
> “Evaluate performance enhancement of Parallel Quicksort using MPI”


 ✅ 🧠 Concept (Short)

 Quicksort → Divide array using a pivot
 Parallel Quicksort (MPI) → split data across processes and sort concurrently
 Uses MPI via mpi4py

---

 💻 ✅ Python MPI Code (Parallel Quicksort + Timing)

> ⚠️ Run using: `mpiexec -n 4 python file.py`

```python id="mpi_qs"
from mpi4py import MPI
import numpy as np
import time

comm = MPI.COMM_WORLD
rank = comm.Get_rank()
size = comm.Get_size()

 ----------- GENERATE DATA (only root) -----------
n = int(input("Enter size of array: ")) if rank == 0 else None
n = comm.bcast(n, root=0)

if rank == 0:
    data = np.random.randint(1, 100, n)
    print("Original Array:", data)
else:
    data = None

 ----------- SCATTER DATA -----------
chunk = np.zeros(n // size, dtype=int)
comm.Scatter(data, chunk, root=0)

 ----------- LOCAL QUICK SORT -----------
def quicksort(arr):
    if len(arr) <= 1:
        return arr
    pivot = arr[len(arr)//2]
    left = arr[arr < pivot]
    middle = arr[arr == pivot]
    right = arr[arr > pivot]
    return np.concatenate((quicksort(left), middle, quicksort(right)))

start = MPI.Wtime()

local_sorted = quicksort(chunk)

 ----------- GATHER RESULTS -----------
gathered = None
if rank == 0:
    gathered = np.zeros(n, dtype=int)

comm.Gather(local_sorted, gathered, root=0)

end = MPI.Wtime()

 ----------- FINAL MERGE (ROOT) -----------
if rank == 0:
    final_sorted = quicksort(gathered)
    print("\nParallel Sorted:", final_sorted)
    print("Parallel Time:", end - start)


 ----------- SEQUENTIAL COMPARISON -----------
if rank == 0:
    start = time.time()
    seq_sorted = quicksort(data)
    end = time.time()
    print("\nSequential Sorted:", seq_sorted)
    print("Sequential Time:", end - start)
```

---

 ⚙️ How it Works (Step-by-Step)

1. Root process generates array
2. Data is divided (Scatter) among processes
3. Each process runs local quicksort
4. Results are collected (Gather)
5. Final merge on root

---

 📊 Performance Idea

| Type           | Time   |
| -------------- | ------ |
| Sequential     | Higher |
| Parallel (MPI) | Lower  |

👉 Speed improves as number of processes increases

---

 🔢 Example

 Input:

```text
[5, 3, 8, 4, 2, 7, 1, 6]
```

 Parallel Distribution:

```text
P1 → [5,3]
P2 → [8,4]
P3 → [2,7]
P4 → [1,6]
```

 Final Output:

```text
[1,2,3,4,5,6,7,8]
```


 🎯 Viva (Oral) Questions & Answers

 🔹 Basics

Q1. What is Quicksort?
Divide-and-conquer sorting algorithm.

Q2. Complexity?
Average: O(n log n)

---

 🔹 MPI Concepts

Q3. What is MPI?
Message Passing Interface for distributed computing.

Q4. What is Scatter?
Distributes data to processes.

Q5. What is Gather?
Collects results from processes.

---

 🔹 Parallel Quicksort

Q6. How is quicksort parallelized?
Data split across processes, sorted independently.

Q7. Why is it faster?
Multiple processors work simultaneously.

---

 🔹 Performance

Q8. What is speedup?
Sequential time / Parallel time.

Q9. What is efficiency?
Speedup / number of processors.

---

 🔹 Limitations

Q10. Overhead in MPI?
Communication cost between processes.


